% The integrated commentary of the concise study, beginning on physical page
% 5. It is built on questions that more than one of the three interpretations
% answers: what is small in the mustard seed and what makes it great; what the
% leaven and the measures are; how far and how long the leavening goes; which
% verse of the Epistle shows the parable in a church; why deliverance is sung
% as done; whose voice the chants are; why a promise must be asked for, with
% the Communion as each interpretation reads it; and the offering. Under each,
% the authorities of the different interpretations are set side by side at the
% text they share. A short comparison of the four senses closes it. Every
% claim, attribution and qualification is the expansive study's; the
% connecting prose is written for this companion.

\runninghead{Detailed Commentary}
\section{The Propers: Detailed Commentary}
\zlabel{triptych:concise:commentary:start}

\subsection{The least of all seeds: what is small, and what makes it great}

Hilary begins from a property of the seed: ``The Lord compared himself to a
grain of mustard, the sharpest and the least of all seeds, whose strength and
power is kindled by tribulations and pressures'' (\work{In Matthaeum}~13.4,
PL~9 col.~994). The seed is sown when the Lord, ``seized by the people and
handed over to death,'' is buried. The herbs it outgrows are the prophets,
whose preaching had been given to a sick Israel as vegetables are given to the
sick; the branches are the apostles; the birds are the nations, tossed about
by the devil's gusts, who ``will fly to them in hope of life.'' Hilary states
the seed's smallness without giving a cause for it.

Ambrose, expounding the same parable in Luke, starts from the other saying in
which mustard appears, faith ``as a grain of mustard seed'' that can move a
mountain. Since both the kingdom and faith are compared to the seed, he
concludes \latin{fides utique regnum coelorum est}, ``faith, then, is the
kingdom of heaven'' (\work{Expositio in Lucam}~VII.176--177). Then he asks what
the seed's nature teaches: ``Its grain is a cheap and simple thing: if it
begins to be ground, it pours out its strength. Faith too at first seems
simple: but if it is ground by adversity, it pours out the grace of its
power'' (VII.178). His example is from his own church, ``our martyrs,'' Felix,
Nabor and Victor, who ``had the scent of faith, but it was hidden'' until,
``crushed by the sword, they spread the grace of their martyrdom to the ends of
the whole world.'' Then he turns to the Lord, who ``was both taken and buried
in a garden; in a garden he grew, where he also rose, and became a tree''
(VII.179). The hearer is to ``sow Christ in your own garden'' (VII.180), and
\latin{Premitur fides, cum crucifixum credimus Christum}, ``faith is pressed
when we believe in Christ crucified,'' as Paul pressed it when he came to the
Corinthians knowing nothing ``but Jesus Christ, and him crucified'' (1~Cor
2:1--2; VII.181). Christ is least, Ambrose says, in appearance and not in
nature, for he ``had no form nor comeliness'' (Is 53:2; VII.183). St Peter
Chrysologus, preaching on Luke's parable, gives the reason from the seed: as
the mustard grain's power is hidden while it is whole and shows most strongly
when it is ground, ``so Christ willed to be ground in his body, who did not
will that his power should stay hidden'' (\work{Sermo}~98).

Jerome's answer is the nearest to these and is still different. The seed is
the preaching of the gospel, ``the least of all teachings,'' because at first
it seems untrue: it preaches ``a man who is God, God dead, and the scandal of
the cross.'' Beside the philosophers' eloquence it looks small, but their
doctrines grow quickly into herbs that wither, while the preaching grows into
a tree (\work{In Matthaeum}~II, PL~26 col.~90). For Jerome the seed is small
because of what it says, and he does not say that the preaching gives its
strength when it is persecuted. Ambrose's ``faith is pressed when we believe
in Christ crucified'' touches Jerome's point, since for both the cross is what
makes the seed look small; but Ambrose also speaks of martyrs ground by the
sword, and Jerome does not. Aquinas follows Jerome: the doctrine seemed least because
it preached God ``suffering, crucified,'' and grew greater than others ``in
solidity, in reach and in usefulness'' (\work{Super Evangelium Matthaei}
c.~13). Augustine gives another reason the seed outgrows the herbs, which are
the opinions of the sects: \latin{ob fervorem fidei}, ``because of the fervour
of faith'' (\work{Quaestiones Evangeliorum}~I.11).

Chrysostom reads the parable as an answer to the disciples' fear. In the
parable of the sower three parts of the seed are lost, and he imagines them
asking, ``And who, and how many will be the faithful?'' The mustard seed
answers that ``in any case the gospel shall be spread abroad'': ``His
disciples were weakest of all, and least of all; but nevertheless, because of
the great power that was in them, It hath been unfolded in every part of the
world'' (\work{Hom.\ in Mt.}~46.2). For Chrysostom, as for Hilary and Ambrose,
the power that makes the small thing great is God's. What is small is, for
him, the disciples; for Hilary the Lord himself; for Ambrose the
believer's faith and the Lord's outward appearance; for Jerome the preaching
of the cross.

\subsection{The leaven: corruption, love or teaching}

Aquinas answers the difficulty that leaven raises with a distinction: the same
thing may be taken \latin{in malum} or \latin{in bonum}, and leaven is evil by
its corruption and good by its heat and its power of spreading through the
dough (\work{Super Evangelium Matthaei} c.~13). Ambrose answers from the very
texts that raise the difficulty. The Lord's warning against the leaven of the
Pharisees and Paul's leaven of malice show that leaven means teaching; and
``the leaven of the cockle is one thing, the leaven of the wheat another''
(VII.190).

Ambrose admits that the parable ``spreads so much ambiguity that it has the
different opinions of many'' (VII.187). He reports that many take the leaven
for Christ, hidden by the Church in the inner parts of the mind, and that some
take the parable of the world; from Matthew's three measures it seemed fitting
to him also to believe the Son of God hidden in the Law, covered in the
Prophets and fulfilled in the Gospel (VII.187--189). Then he gives his own
view: ``Yet I am glad to follow what the Lord himself taught, that the leaven
of the Church is spiritual teaching'' (VII.190). Jerome reaches a similar
answer more briefly: the woman ``seems to me'' to be the apostolic preaching or
the Church, and the leaven the knowledge of the Scriptures (PL~26 col.~91).
Augustine gives the other answer in two lines: ``by the woman he means wisdom;
by the leaven, love, which warms and rouses'' (\work{Quaestiones
Evangeliorum}~I.12). Bede, on Luke, takes up Augustine's words and adds their
object, love ``which warms and rouses the mind'' (\work{In Lucam}~IV, on
13:21).

On the three measures these four agree more closely than on the leaven: the
measures are the whole person. Augustine names them from the great
commandment, ``with the whole heart, and with the whole soul, and with the
whole mind'' (Mt 22:37). Bede makes the commandment the reason the leaven is
hidden in three measures at all: ``The woman hides the leaven of love in three
measures of meal, because the Church commands us to love the Lord with the
whole heart, the whole soul and the whole strength.'' Jerome takes the
measures as spirit, soul and body, hidden with the Scriptures ``so that spirit,
soul and body, brought into one, may not be at odds among themselves''; after
Plato's three parts of the soul he adds, of himself and his readers, ``that in
reason we may possess prudence; in anger, hatred of vices; in desire, longing
for the virtues'' (PL~26 col.~91). Ambrose takes them as flesh, soul and
spirit, the reborn man whom the Church sanctifies with spiritual leaven
(VII.190--191). Bede says also what the meal is. The woman is the Church, and
we are her meal, ``as many of us as are ground by the exercise of fear and
hope, as between an upper and a lower millstone, so that \dots\ we being many
may be one bread, one body in Christ'' (1~Cor 10:17).

In Bede's sentence the Church hides love in the three measures by commanding
it: a teaching, the great commandment, is the means, and love is what is
hidden. The second interpretation takes the two answers, teaching and love,
in that order: the teaching is received first, and the love it commands works
through the whole person.

Others answer differently again. Hilary takes the leaven for the Lord, hidden
by the synagogue ``through the sentence of death,'' and the three measures for
the Law, the Prophets and the Gospels, made one by it (\work{In
Matthaeum}~13.5); Chrysologus too reads the measures so. Over a Trinitarian
reading of the measures the Fathers truly disagree. Jerome reports it as
others' and calls it ``a pious sense indeed, but the understanding of a parable
and the doubtful meaning of riddles can never advance to the authority of
doctrines''; Hilary writes, ``I do not know whether reason permits us to think
so,'' since Father, Son and Holy Ghost need no leaven mixed in from outside
(13.6); Chrysologus preaches it from Sara's three measures (\work{Sermo}~99).

The strongest alternative is Chrysostom's. For him the leaven is the
disciples, mixed into the multitude of the world, and the three measures mean
only ``many'': ``even so shall ye convert the whole world'' (\work{Hom.\ in
Mt.}~46.2). Aquinas reports this reading with the apostles hidden ``in
tribulations.'' It makes the leaven a second parable of the mission, and so it
belongs with the first interpretation, beside the mustard seed.

\subsection{Until the whole is leavened: how far, and when}

The parable's last words, \latin{donec fermentátum est totum}, raise two
questions: how much of the meal the leaven changes, and when it has finished.
Chrysostom, though he takes the leaven for the disciples, explains the
parable's verb: ``He said not, `put,' simply, but `hid','' so that the leaven
is ``actually mixed''; and ``though it be buried, yet is not destroyed, but by
little and little transmutes all into its own condition'' (46.2). Ambrose calls
the Church's leavening ``not hasty, nor haphazard, but brought about by long
meditation,'' so that the three may be one; the woman goes on mixing ``until
the whole inner man of the heart, hidden, is leavened, and rises into the
grace of heavenly bread,'' and the teaching is Christ's leaven ``because Christ
is bread'' (VII.191).

Bede argues from \latin{donec}. The leaven was hidden until the whole was
leavened ``because love hidden in our mind must grow until it changes the
whole mind into its own perfection, so that the soul can love, do and recall
nothing but the love of its maker.'' In his sentence the parable's ``until''
becomes a requirement, and he places its end beyond this life: ``This is
begun here, but it is to be perfected there, where, God being all in all, he
warms all with one and the same fire of his love'' (1~Cor 15:28; \work{In
Lucam}~IV, on 13:21). Chrysologus asks the question the same word raises:
\latin{fermentatur totum. Quando? Adveniente Christo}, ``the whole is
leavened. When? When Christ comes'' (\work{Sermo}~99).

The second and third interpretations use these answers for different ends. In
the second, Bede's ``must grow'' is a demand on the person, and love is
finished in no one in this life. In the third, Chrysologus's ``when Christ
comes'' puts the parable's end at the Lord's coming: the seed is sown and the
leaven hidden, while the tree and the whole leavened are still to come. The
Postcommunion is read in all three. Those who pray it have been fed,
\latin{pasti}, and ask that they may \latin{semper}, always, long for
\latin{éadem}, the same things, ``by which we truly live.'' On Bede's account
of love that continuing desire is the leaven's work going on in those who
have received. In the third interpretation the meal has been given and the
longing continues, as in the Epistle Jesus has delivered and the
Thessalonians wait. In the first, Bede's remark that the preaching grows also
``in the mind of each believer, because no one is made perfect all at once''
(\work{In Lucam}~IV, on 13:19) means that those who have been fed are still
growing.

\subsection{Three verses of the Epistle}

Each interpretation takes a different verse of Paul's thanksgiving. The first
takes verses 6 and 8. The Thessalonians received the word \latin{in
tribulatióne multa}, and two verses later the word of the Lord was spread
abroad from them \latin{in omni loco}, in every place. Paul tells of the young church in the order
of Ambrose's martyrs, crushed first and then heard of everywhere. Chrysostom
says why the pressure came at once: ``God permitted that persecutions should
arise immediately, that no one might say, that the Preaching was established
lightly or by flattery''; and of the work of faith, ``If thou believest, suffer
all things; if thou dost not suffer, thou dost not believe'' (\work{Hom.\ in
1~Thess.}~1). Aquinas, on the Latin, says that it is the joy of the Holy
Spirit and no other joy, the Spirit ``who is the love of God, who gives joy to
those who suffer for Christ, because they love him'' (\latin{qui est amor Dei,
qui facit gaudium patientibus propter Christum, quia amant eum}; \work{Super I
ad Thessalonicenses} c.~1 l.~1). On verse~5, where the gospel came ``in power
also,'' he gives two explanations: the preaching was confirmed ``not by
arguments but by the power of miracles,'' or Paul showed them ``a virtuous
life.'' Paul had reasoned in their synagogue that ``the Christ was to
suffer'' (Acts 17:3), and Ambrose says that faith is pressed when we believe
in Christ crucified. By Ambrose's account, then, the Thessalonians' faith
was pressed both in what they believed and by the persecution that followed.

The second takes verse~3, \latin{óperis fidei vestræ, et labóris, et
caritátis}. The Greek that Chrysostom expounds reads ``the labour of love,''
and he asks what labour it is to love: ``Merely to love is no labor at all. But
to love genuinely is great labor,'' when a thousand things would draw us from
love and we hold out against them all. His example is Jason, dragged before
the rulers at Thessalonica for having received Paul (\work{Hom.\ in
1~Thess.}~1). Theodoret finds the labour in bearing a brother's faults, his
envy, anger, pride and ingratitude (\work{Interpretatio in 1~Thess.}, on
1:3). Aquinas, reading the Missal's Latin, attaches the labour to faith:
faith is not enough ``unless it has work and labour,'' and then he names
charity. If the leaven is love, as Augustine and Bede say, the charity
Paul names is that leaven, and the work and the labour are where it shows.

The third takes the last verse, the waiting for the Son. Chrysostom names what
is waited for, ``the Resurrection, and the Ascension, and the second Coming,
the Judgment,'' and the condition of those who wait: ``The terrible things are
in hand, but the good things are in the future \dots\ See how much hope is
required'' (\work{Hom.}~2). His translator's ``which delivereth us'' follows
the Greek; the Missal's Latin is the perfect \latin{erípuit}. Paul also
remembers the enduring of the hope, and Aquinas says that hope ``makes us bear adversity patiently,''
``before God, not before the eyes of men.''

The Collect's pair, word and deed, is read beside these verses differently in
each. Paul's pair in verse~5 is word and power. Chrysostom says that it was
not the apostles' signs that made them admirable but ``their contempt of
wealth, their despising glory \dots\ it is their life, so bright on all sides,
which also draws down the grace of the Spirit'' (\work{Hom.\ in Mt.}~46.3). On
Aquinas's second explanation the power of verse~5 is a virtuous life, close
to the life Chrysostom praises, and in the first interpretation the deeds the
Collect asks for can be understood as a life of that kind. In the second, the Collect's people meditate
\latin{semper}, always; their constant thought may be taken as their share in
the slow work that Ambrose says the Church does ``by long meditation,'' and the
\latin{et \dots\ et} as asking that it reach their speech and their action. In
the third, the same \latin{semper} covers the whole time of waiting.

\subsection{Deliverance sung as done}

The third interpretation begins from the tenses set out above, and Augustine
meets the question they raise on the Gradual's own verses. ``Thou
hast saved us from our enemies'' is ``spoken of the future under the figure of
the past. But this is the reason that it is spoken of as if it were past, that
it is as certain as if it were past.'' His precedent is the Passion, foretold
in the psalm as ``They pierced My hands and My feet'': ``to God things to come
also are as certain as if they were past.'' On the next verse he observes
``how he intermingles words expressive of a future time'' (\work{En.\ in
Ps.}~43.6--7). The Gradual sings exactly that pair of verses. On the Introit's
psalm he asks how the psalmist can pray, ``Convert us, O God,'' for what he
has just called done, ``Thou hast forgiven the iniquity of thy people'' (Ps
84:3, 5), and answers in the psalmist's voice: ``These things I speak of as
done, because I see them about to be done: but because they are not yet done,
I pray that they may come, which I have already seen'' (84.4).

Aquinas answers from the cause. On the Epistle's \latin{erípuit} he says that
we still await deliverance from the punishment to come, but ``From the cause
of the punishment, namely from sin, we are freed through Christ. Therefore he
says, who delivered us'' (\work{Super I ad Thessalonicenses} c.~1 l.~1). On
the psalm's \latin{Salvasti}, the Vulgate's verb where the chant has
\latin{Liberásti}, he gives both times, the present, ``because
though they harm us in body they cannot harm us in spirit,'' or the future
(\work{Super Psalmos}~43, n.~3). Augustine stresses the certainty of what God
has foretold, Aquinas the removal of the cause now and of the penalty later,
Chrysostom the hope that present terrors require.

The plainest objection gives the past tenses a past. Psalm~43 remembers how
God planted his people in the land, and Jeremias's promise was made to
particular exiles and kept when they came back from Babylon. Jerome grants
that history: he reads the letter first as a letter to captives, and only
then of the Church. On Mt 13:35 he states his rule. Because the psalm Matthew
quotes is spoken ``in the person of the Lord'' and narrates the signs of the
exodus, ``all those things that are written are to be taken as parables: the
letter does not sound only what is plain, but hidden mysteries as well''
(\work{In Matthaeum}~II, PL~26 col.~93). By that rule Jeremias's letter is
history and also more. Augustine, on the same psalm, says that ``the mystery
of the Kingdom of Heaven was veiled in the Old Testament, which in the fulness
of time should be unveiled in the New'' (\work{En.\ in Ps.}~77.1--2); and he
hears Psalm~43 as the voice of a persecuted Church because Paul quoted it to
one, so that its ``Thou hast saved us'' is a deliverance still to come.

Three medieval commentators on the Mass who expound these chants read them as
one sequence. Rupert of Deutz takes the Introit's thoughts of peace as the promise of the
banquet of grace to Christ's brethren according to the flesh, prefigured in
Joseph, who received the brothers who had sold him: in the Gradual they give thanks for release ``from the captivity of souls,''
in the Offertory they ask the true Joseph to forget their wickedness (Gen
50:17), and the Communion is his answer, ``Fear not'' (\work{De divinis
officiis}~XII.23); William Durandus takes up Rupert's Joseph
(\work{Rationale}~VI.140). Sicard of Cremona reads the same four chants as the
remnant of Israel answering the promise, and in the Communion the Lord ``gives
assurance of salvation'' (\work{Mitrale}~VIII.23). For both Rupert and Sicard
the sequence belongs to the remnant of Israel at the end of time, and they put
the Gradual's thanks in that future although \latin{Liberásti} is a perfect,
as is the Offertory's \latin{clamávi}. Like Augustine they take the
Gradual's past tense of something still to come, though neither argues from
the tense as he does.

\subsection{Whose voice the chants are}

The first interpretation hears the Gradual and the cry from the depths as the
voices of the afflicted. The Introit sings God's words to a small and captive
people far from their land, and Jerome's gloss on the clause the antiphon
omits speaks of ``patience in the present labours.'' Augustine hears the
martyrs in Psalm~43. Chrysostom takes the Gradual's ``all the day'' as the
whole of life, in which we glory in God not only when he helps us but also
when he forsakes us, as Paul glories only in the cross (Gal 6:14;
\work{Expositio in Ps.}~43). Those who are pressed cry from the second of
Hilary's three depths, ``the deep surge of human calamities,'' and Chrysostom
says that prayer from the deepest heart, as the souls of the afflicted pray,
stands like a deep-rooted tree against every gust of the devil, while prayer
from the lips is overturned (\work{Expositio in Ps.}~129.1).

For the second, the cry is prayer that changes the one who prays. Chrysostom
says on the same verse that such prayer, even before it receives what it asks,
calms the passions, softens anger and drives out envy. The second
interpretation likens that inward change to the leaven's work in the meal,
hidden and from within. Aquinas says on
the Gradual that the saints are afflicted in body by their enemies and in
spirit ``when those who live among them sin'' (\work{Super Psalmos}~43,
n.~3), and bearing such neighbours comes close to the labour Theodoret assigns
to love.

The third hears in it the prayer of those who are still in the depths and
waiting. Hilary asks why we must pray at all to a God who is everywhere and
hears everything, and answers: \latin{Naturae suae est ut audiat; sed fidei
nostrae est ut precemur audiri}, ``It is of his nature that he hears; but it
is of our faith that we pray to be heard'' (\work{Tractatus in
Ps.}~129.2--3).

\subsection{Why a promise must be asked for}

Jerome asks a question close to Hilary's about the Introit's exiles. God has
promised to bring them back; why must they pray? ``Surely the Lord could have
fulfilled what he had promised without the calling and prayer of the
captives; but he urges them to prayer, so that they may deserve to receive
what was promised.'' He quotes \latin{Petite, et accipietis}, ``Ask, and you
shall receive'' (Mt 7:7; \work{In Hieremiam}~V, PL~24 col.~893). He reads the
letter anagogically of the Church, in an order: we remain in the confusion of
this world until we deserve the rest of the number seven, and then,
repentance having been received, God fulfils what he promised and brings us
back ``to our place, the Church.'' He continues, ``he will give us the end of
our labour and patience: and we shall call upon him, and return to the
Church, and pray, and be heard'' (col.~894), which is his reading first of the
clause the antiphon omits and then of the antiphon's own words. For Jerome,
prayer makes the exiles worthy to receive what is promised; for Hilary,
praying to be heard is the part faith plays.

The Communion's promise, \latin{crédite quia accipiétis}, is conditioned
differently in each interpretation. In the first, it is spoken to
mountain-moving faith. Ambrose began his exposition of the mustard seed from
the saying that such faith moves a mountain, and Jerome, on the parable, says
that ``what the Lord said is done by faith that is compared to a grain of
mustard, the Apostle teaches can be done by all faith'' (1~Cor 13:2; PL~26
col.~91): mustard-seed faith, for both, is great faith. In Mark the
Communion's verse follows directly on the mountain, so the promise is spoken
to the faith both Fathers call the mustard seed. Bede, on Mark, says that ``he
who has perfect faith, which works through love, can by praying, or even by
commanding, move spiritual mountains,'' and gives Paul's examples, among them
the spirit cast out of the girl at Philippi, ``a most proud mountain,'' whose
casting into the sea brought persecution on Paul at once (\work{In
Marcum}~III, on 11:22--25). Paul recalls that persecution to the
Thessalonians (1~Thess 2:2). Chrysostom, on Matthew's form of the promise (Mt
21:21--22), says it was given ``so that they might not be afraid and tremble
at plots against them'' (\work{Hom.\ in Mt.}~67.2).

In the second interpretation the promise is addressed to the whole person
made one. Jerome's
sentence on the leaven continues: spirit, soul and body brought into one, ``but
when two and three have agreed, may obtain from the Father whatever they ask''
(Mt 18:19; PL~26 col.~91). Ambrose puts the full obtaining at the
resurrection, when there will remain ``an incorrupt communion of body and soul
and spirit, and then we shall obtain all that we ask'' (VII.192). Neither
comments on Mark, and Mt 18:19 is another saying; if prayer obtains when the
whole person agrees, as both say, the Communion's ``whatsoever you ask'' is
spoken to such a person.

For the third, the promise is spoken to prayer itself. The antiphon commands in
the present and promises in the future, the form of \latin{Petite, et
accipietis} which Jerome quotes for the exiles, and with Hilary's sentence in
mind its \latin{crédite} can be heard as a call to the faith that does the
asking. Bede adds a condition, which the second interpretation also uses, for
those whose faith is not yet perfect: they are to ask that their sins be
forgiven, ``and they will without doubt obtain what they ask, provided that
they first forgive others who sin against them'' (\work{In Marcum}~III, on
11:25). Mark's next verse gives the command to
forgive.

\subsection{The offering}

The Church holds that the offering of the Mass is the sacrifice of the cross
made present. In the first interpretation that is the death in which Hilary
sees the seed sown, and the Secret asks that this offering cleanse and renew,
govern and protect, those who offer it. In the second, the Secret asks for cleansing
and renewal before governing and protection, and Paul's command about leaven
also puts purging before newness: ``Purge out the old leaven, that you may be
a new paste'' (1~Cor 5:7), the verse Aquinas cites for leaven taken
\latin{in malum}. Read beside it, the cleansing the Secret asks for is the
purging of the old leaven. The keeping it asks for, \latin{gubérnet et
prótegat}, is in the third interpretation asked for the time of waiting.

\subsection{The four senses compared}

Literally, the first interpretation starts from the seed that becomes a tree
and a church persecuted at once whose faith is heard of everywhere; the
second, from leaven worked through a large batch of meal and a church whose
faith works, labours and loves; the third, from Jeremias's letter, Paul's
converts waiting for the Son, and Jesus speaking in parables of things hidden
from the foundation of the world. The allegory of the first is Christ seized
and buried (Hilary, Ambrose), risen in the garden and become a tree
(Ambrose), ground in his body so that his power should not stay hidden
(Chrysologus), with the apostles for branches (Hilary); of the second, the
Church mixing Christ's teaching into
the reborn man until he rises ``into the grace of heavenly bread'' (Ambrose),
and ourselves her meal, ground into one bread (Bede); of the third, the
captivity of sin turned away by Christ (Augustine) and the Church as the
exiles' place of return (Jerome). Morally, the first asks faith that pours
out its strength when adversity grinds it (Ambrose) and, believing, suffers
all things (Chrysostom); the second, love that changes the whole mind (Bede),
love that is labour (Chrysostom) and labours in bearing a brother's faults
(Theodoret), with reason, anger and desire made prudence, hatred of vice and
longing for virtue (Jerome); the third, prayer for what God has promised
(Jerome), the faith that prays to be heard (Hilary), and, for those not yet
perfect, forgiving others so as to obtain the forgiveness they ask (Bede). The
anagogy of the first is the tree into whose branches the nations fly and
flourish ``in the security of a sure station'' (Ambrose, VII.185--186); of
the second, love perfected where
God is all in all (Bede) and ``an incorrupt communion'' for those who love
Christ (Ambrose, VII.194); of the third, the whole leavened at the Son's coming
(Chrysologus) and the return after the confusion of this world, when we shall
call and be heard (Jerome). The first places the kingdom's hiddenness in
suffering and its growth outward to the nations, the second in the person and
its work inward, the third in time and its fulfilment at the Son's coming. The
real disagreements among the Fathers cited here lie within the interpretations,
over what the seed, the leaven and the measures are.
